\documentclass[11pt]{article}
\usepackage[margin=1in]{geometry}
\usepackage{amsmath,amssymb}
\usepackage{hyperref}
\usepackage{microtype}
\newcommand{\ii}{\mathrm{i}}

\title{The Long Wall in a Time-Symmetric Two-Boundary Model}
\author{Mike Stay}
\date{October 2026}

\begin{document}
\maketitle

\begin{abstract}
A useful objection to source-driven pilot-wave models comes from a modified double-slit experiment in which a long splitter plate is inserted upstream of the slits.  Once a localized particle chooses one channel, the wave in the other channel is no longer driven.  If that branch disperses, its relative amplitude tends to zero and so does the field-interference visibility.  Andersen and collaborators proposed this as a test of wave-driven droplet models.

We examine a narrower question for a time-symmetric construction.  In a scalar paraxial point-slit model, condition on an advanced contribution associated with a candidate detector event.  Its two slit amplitudes have the standard double-slit relative phase, with no dependence on the splitter length.  If subsequent backward propagation through the two matched channels acts on both branches by the same nonzero scalar transfer function, it changes neither their amplitude ratio nor their relative phase.

This is a conditional common-mode result, not a complete solution of the long-wall experiment.  We do not derive the Green function of a physical splitter, the strength of a self-consistent absorber response, the visibility after the advanced contribution is combined with other field components, or the probability distribution of detector events.  The calculation addresses one propagation component of the long-wall objection and identifies the further assumptions that a complete time-symmetric account would require.
\end{abstract}

\section{Introduction}

The double-slit experiment is a good place to test stories about quantum mechanics, because the experimental facts are simple while the stories are not.  A particle is detected at one place on a screen.  Nevertheless, changing the geometry of a slit through which the particle did not pass changes the statistics of where particles arrive.

A particularly sharp version of this problem was proposed by Andersen, Madsen, Reichelt, Ahl, Lautrup, Ellegaard, Levinsen and Bohr in their study of wave-driven droplets \cite{Andersen2015}, and emphasized by Bohr, Andersen and Lautrup \cite{Bohr2016}.  Insert a long splitter plate upstream of a double slit.  A classical particle must choose one side of the plate near its upstream end.  If the guiding wave is continually generated by the particle, the wave in the chosen channel remains driven, while the wave in the other channel is left to disperse.  Their source-driven model therefore predicts that a sufficiently long plate makes the empty-channel field too weak to produce the usual two-path effect.  The cited papers present the argument and its predicted behavior; they do not report an experimental sweep of splitter length.

This is a powerful objection to a large class of mechanical analogues of quantum mechanics.  But it combines several logically distinct questions.  Does the field reach both slits?  Do the two branches acquire an additional relative phase or differential attenuation in the long channels?  How strong is the total field relative to a particle-attached branch?  Finally, what rule converts a field into detector-event probabilities?

There are theories in which the physical solution is constrained by data at both ends of an experiment.  Wheeler and Feynman formulated electrodynamics using retarded and advanced interactions \cite{WheelerFeynman1945,WheelerFeynman1949}.  Cramer's transactional interpretation describes a quantum event using an offer wave and an absorber response \cite{Cramer1986}.  Wharton and others have studied formulations in which initial and final boundary conditions jointly constrain the solution \cite{Wharton2007}.  Hasselmann's metron program is another example: its published diffraction treatment is largely forward-propagating, while the broader theory explicitly appeals to microscopic time-reversal symmetry \cite{Hasselmann1997,Hasselmann1998,Hasselmann2004}.

The modest question addressed here is:

\begin{quote}
In an effective two-channel model, does extending a symmetric common channel necessarily introduce an additional relative phase or differential attenuation between two branches conditioned by a future boundary?
\end{quote}

Under a precise scalar-transfer hypothesis, the answer is no.  This does not settle the complete dynamical long-wall objection.  It isolates the part that is simply a question of relative propagation.

\section{The forward-only visibility argument}

Consider a double slit in the plane $x=0$, with slit centers
\begin{equation}
 y_\pm=\pm\frac{a}{2}.
\end{equation}
A detector is a distance $d$ downstream, and a splitter plate extends upstream to a tip at $x=-L$.

Suppose the occupied slit is illuminated with amplitude $A$, while the empty slit is illuminated with amplitude $r(L)A$.  Let $f(\theta)$ be their common single-slit envelope and let $\delta(\theta)$ be their relative phase.  After choosing a symmetric phase convention, the coherent field may be written
\begin{equation}
 \Psi(\theta)
 =A f(\theta)
 \left[
 e^{\ii\delta(\theta)/2}
 +r(L)e^{-\ii\delta(\theta)/2}
 \right].
 \label{eq:forwardfield}
\end{equation}
This is the standard unequal-amplitude two-wave parametrization; see, for example, \cite{Daniel2022}.

Taking the modulus squared gives
\begin{align}
 |\Psi(\theta)|^2
 &=|A f(\theta)|^2
 \left[
 1+r(L)^2+2r(L)\cos\delta(\theta)
 \right].
 \label{eq:forwardintensity}
\end{align}
When $\delta$ is varied through a full fringe, the maximum and minimum are proportional to $(1+r)^2$ and $(1-r)^2$.  The coherent field-intensity visibility is therefore
\begin{equation}
 \mathcal V(L)
 =\frac{I_{\max}-I_{\min}}{I_{\max}+I_{\min}}
 =\frac{2r(L)}{1+r(L)^2}.
 \label{eq:visibility}
\end{equation}
Thus
\begin{equation}
 r(L)\longrightarrow0
 \quad\Longrightarrow\quad
 \mathcal V(L)\longrightarrow0.
\end{equation}

The algebra is uncontroversial.  The physical premise is that the empty-channel branch is populated only by a forward-propagating field which becomes negligible as the common wall is lengthened.  Moreover, \eqref{eq:visibility} is a statement about coherent field-intensity contrast.  Applying it directly to particle arrivals requires an additional dynamical or detector-response law.

\section{A conditioned two-boundary point-slit calculation}

We now condition on a candidate detector interaction at transverse coordinate $Y$ in the plane $x=d$.  The calculation is deliberately simple: a monochromatic scalar field, point evaluation at the slit centers, and free paraxial propagation between the detector and slit plane.

The one-dimensional paraxial kernel over longitudinal distance $\ell>0$ is
\begin{equation}
 K_\ell(y_2,y_1)
 =
 \left(\frac{k}{2\pi\ii\ell}\right)^{1/2}
 \exp\left[
 \frac{\ii k(y_2-y_1)^2}{2\ell}
 \right],
 \label{eq:kernel}
\end{equation}
where $k=2\pi/\lambda$.  The paraxial generator $H=-\partial_y^2/(2k)$ is self-adjoint on its standard domain, so the forward evolution is unitary and backward propagation is represented by the adjoint kernel $K^*$.

Let $A_D$ denote the amplitude assigned to the conditioned future interaction.  This amplitude is an input to the toy boundary-value calculation; it is not derived here from an absorber model.  Backward propagation to the two slit centers gives
\begin{equation}
 A_\pm(Y)
 =A_DK_d^*\left(Y,\pm\frac a2\right).
 \label{eq:slitamplitudes}
\end{equation}
In the paraxial point kernel, the prefactor is independent of transverse separation, so
\begin{equation}
 |A_+(Y)|=|A_-(Y)|.
 \label{eq:equalmagnitudes}
\end{equation}
Their ratio is
\begin{align}
 \frac{A_-(Y)}{A_+(Y)}
 &=\exp\left\{-\frac{\ii k}{2d}
 \left[
 \left(Y+\frac a2\right)^2
 -\left(Y-\frac a2\right)^2
 \right]\right\} \\
 &=\exp\left(-\frac{\ii kaY}{d}\right).
 \label{eq:relativephase}
\end{align}
Thus an unwrapped relative phase is
\begin{equation}
 \Delta\phi(Y)=-\frac{kaY}{d},
 \qquad\text{understood modulo }2\pi.
\end{equation}
No splitter length $L$ occurs in this detector-to-slit calculation.

The coherent sum factorizes as
\begin{equation}
 A_+(Y)+A_-(Y)
 \propto
 2\exp\left[-\frac{\ii k}{2d}
 \left(Y^2+\frac{a^2}{4}\right)\right]
 \cos\left(\frac{kaY}{2d}\right).
 \label{eq:conditionedsum}
\end{equation}
Consequently the conditioned two-path amplitude has modulation
\begin{equation}
 \mathcal M(Y)
 \equiv|A_+(Y)+A_-(Y)|^2
 \propto4\cos^2\left(\frac{kaY}{2d}\right),
 \label{eq:modulation}
\end{equation}
with spacing
\begin{equation}
 \Delta Y=\frac{\lambda d}{a}.
 \label{eq:spacing}
\end{equation}
Equation \eqref{eq:modulation} is not, by itself, a probability distribution for detector events.  It is the two-path modulation of a coherent amplitude conditioned on a candidate event at $Y$.

The exact three-dimensional free Green function adds a familiar qualification.  If
\begin{equation}
 R_\pm=\sqrt{d^2+\left(Y\mp\frac a2\right)^2},
\end{equation}
then an outgoing spherical Green function has magnitude proportional to $1/R_\pm$.  The two magnitudes are not exactly equal away from $Y=0$.  In the paraxial far field, their difference is small and the path-difference phase reduces to \eqref{eq:relativephase}.  This ordinary detector-to-slit correction depends on $a,d,Y$, not on the upstream wall length $L$.

\section{The effective-channel common-mode result}

The calculation above establishes the relative data at the slit mouths.  To discuss the extra wall length, introduce abstract transfer functions for backward propagation through the two channels:
\begin{equation}
 B_+(L,Y)=T_+(L)A_+(Y),
 \qquad
 B_-(L,Y)=T_-(L)A_-(Y).
 \label{eq:channeltransfer}
\end{equation}
Then
\begin{equation}
 \frac{B_-(L,Y)}{B_+(L,Y)}
 =\frac{T_-(L)}{T_+(L)}
 \frac{A_-(Y)}{A_+(Y)},
 \label{eq:generalratio}
\end{equation}
whenever the denominators are nonzero.

\medskip
\noindent\textbf{Proposition.}
\emph{Suppose the relevant matched mode in each channel undergoes identical nonzero scalar transfer,
\begin{equation}
 T_+(L)=T_-(L)=T(L)\neq0.
 \label{eq:scalartransfer}
\end{equation}
Then extending the common channel length does not change the amplitude ratio or relative phase of the two branches:}
\begin{equation}
 \frac{B_-(L,Y)}{B_+(L,Y)}
 =\frac{A_-(Y)}{A_+(Y)}.
 \label{eq:commoncancel}
\end{equation}

The proof is the cancellation of the common factor $T(L)$.  Notice that this algebra is not intrinsically time-directed: the same statement holds for forward or backward propagation.  Time symmetry matters earlier, by providing a possible reason for a field branch in the nominally empty channel to be constrained by a future boundary.

The proposition is an effective-channel result.  It does not follow from evaluating a free-space Fresnel kernel between the slit centers and the upstream tip.  A physical splitter imposes boundary conditions along the wall and may support finite-width modes, losses, mode conversion, or coupling.  The present model assumes that the two relevant channel modes are matched, uncoupled, and acted on by the same scalar.  Deriving \eqref{eq:scalartransfer} for a specified wall would require solving or characterizing that wall's boundary-value problem.

\section{Why mirror symmetry alone is insufficient}\label{sec:exchange}

The scalar hypothesis \eqref{eq:scalartransfer} is stronger than channel-exchange symmetry.  To see this, let
\begin{equation}
 J=
 \begin{pmatrix}
 0&1\\
 1&0
 \end{pmatrix}
\end{equation}
exchange the two channels, and define the unitary, exchange-symmetric operator
\begin{equation}
 U(\vartheta)=\cos\vartheta\,I+\ii\sin\vartheta\,J.
\end{equation}
It commutes with $J$, but it mixes the channels.  Acting on the amplitude vector $(1,\ii)^T$ gives
\begin{equation}
 U(\vartheta)
 \begin{pmatrix}1\\ \ii\end{pmatrix}
 =
 \begin{pmatrix}
 \cos\vartheta-\sin\vartheta\\
 \ii(\cos\vartheta+\sin\vartheta)
 \end{pmatrix}.
\end{equation}
At $\vartheta=0$ the output magnitudes are equal; at $\vartheta=\pi/4$ the first component vanishes.  Symmetry and losslessness therefore do not guarantee preservation of an arbitrary amplitude ratio.  The proposition requires identical scalar action on the relevant modal subspace, or some equally strong replacement.

\section{Advanced coherence is not total-field visibility}\label{sec:total}

Even perfect common-mode propagation of the advanced pair does not fix the visibility of a total field containing other contributions.  Write
\begin{equation}
 \Phi_+(L,Y)=R_+(L,Y)+t(L)A_+(Y),
 \qquad
 \Phi_-(L,Y)=R_-(L,Y)+t(L)A_-(Y).
 \label{eq:totalfield}
\end{equation}
The factor $t(L)$ cancels from the ratio of the two advanced pieces, but generally not from $\Phi_-/\Phi_+$.

A minimal counterexample is
\begin{equation}
 R_+=1,
 \qquad R_-=0,
 \qquad A_+=A_-=1.
\end{equation}
Then
\begin{equation}
 \Phi_+=1+t,
 \qquad
 \Phi_-=t,
\end{equation}
and the coherent two-path visibility is
\begin{equation}
 \mathcal V_{\rm total}(t)
 =\frac{2|t|\,|1+t|}{|1+t|^2+|t|^2}.
 \label{eq:totalvisibility}
\end{equation}
If $t>0$ is real and tends to zero, then $\mathcal V_{\rm total}\to0$ despite perfect coherence within the advanced pair.  Even fixing $|t|=1$ is insufficient: $t=1$ gives visibility $4/5$, while $t=-1$ cancels the first channel and gives zero visibility.

A simple sufficient condition for an $L$-independent normalized total pattern would be factorization of the entire relevant field,
\begin{equation}
 \Phi_\pm(L,Y)=g(L)b_\pm(Y),
 \qquad g(L)\neq0,
 \label{eq:totalfactorization}
\end{equation}
with $b_\pm$ independent of $L$.  Nothing in the point-slit phase calculation establishes \eqref{eq:totalfactorization}.

\section{Four levels of the long-wall problem}

It is helpful to distinguish four levels which are easily run together.

\begin{enumerate}
\item \textbf{Detector-to-slit relative phase.}  In the stated paraxial point-slit model, this is calculated explicitly and has no $L$ dependence.
\item \textbf{Propagation through the channels.}  The ratio is preserved if the relevant two modes undergo identical nonzero scalar transfer.  This is an assumption of the effective-channel model, not a derived wall Green function.
\item \textbf{Normalization of the complete source--field--absorber solution.}  This determines the absolute advanced response and its interference with retarded or particle-attached pieces.  It is not calculated here.
\item \textbf{Distribution of detector events.}  Conditioning on a candidate $Y$ does not select $Y$ or assign it a probability.  A detector-response or probability law is still required.
\end{enumerate}

The paper settles a conditional result at the first two levels.  It does not settle the third or fourth.

\section{Relation to transactional and metron pictures}

The result is not tied to one interpretation, but each interpretation supplies additional structure which the toy calculation does not reproduce.

In Cramer's transactional interpretation, an absorber returns a confirmation wave in response to the offer wave reaching it \cite{Cramer1986}.  The response normalization is therefore part of the source--absorber dynamics; it is not an independently chosen constant.  In the present calculation $A_D$ is simply a conditioned boundary amplitude.  The calculation checks the relative phase produced by propagating it through the two slits, not the mechanism which determines its magnitude.

In two-boundary formulations, the same observation can be stated without offer and confirmation waves: future boundary data may participate in determining the field in the nominally empty channel \cite{Wharton2007}.  Again, a complete theory must determine which final boundary occurs and with what weight.

Hasselmann's metron model is a useful case study because it combines a localized particle core with an extended de Broglie component.  Its published diffraction treatment is essentially a forward scattering model.  Hasselmann and Hasselmann compared their calculations with quantum-theoretical diffraction predictions and discussed their relation to neutron experiments \cite{Hasselmann2004}.  The broader program also invokes microscopic time-reversal symmetry \cite{Hasselmann1998}.  The present calculation does not complete that program; it identifies one conditional propagation property that a two-boundary completion could use.

Walking droplets have a different mathematical structure: a source-driven wave in a dissipative medium is treated as an initial-value problem.  The long-splitter argument predicts that the undriven branch eventually becomes too weak to affect the droplet \cite{Andersen2015,Bohr2016}.  The common-factor identity does not refute that prediction.  It asks what follows if a separate future-conditioned branch exists and if its two channel transfers are matched.

\section{What has and has not been shown}

The following statements have been established within the stated model:

\begin{itemize}
\item the conditioned detector-to-slit ratio is $\exp(-\ii kaY/d)$ in the paraxial point-slit approximation;
\item the corresponding coherent amplitude has the standard cosine modulation;
\item identical nonzero scalar channel transfers cannot change a two-component amplitude ratio;
\item neither mirror symmetry alone nor coherence of the advanced pair alone guarantees preservation of the total visibility.
\end{itemize}

The following have not been established:

\begin{itemize}
\item the Green function or transfer operator of a specified physical wall;
\item absence of channel mixing or mode conversion in that wall;
\item the absolute magnitude and phase of a self-consistent absorber response;
\item factorization of the complete field as in \eqref{eq:totalfactorization};
\item a guidance, absorption, or detection law converting the conditioned modulation into event frequencies;
\item the Born distribution over candidate detector positions.
\end{itemize}

The algebraic identities in the preceding sections --- the detector-to-slit ratio \eqref{eq:relativephase}, the sum factorization \eqref{eq:conditionedsum}, the common-mode cancellation \eqref{eq:commoncancel}, the $U(\vartheta)$ counterexample of \S\ref{sec:exchange}, the visibility formula \eqref{eq:visibility}, and the total-field counterexamples of \S\ref{sec:total} --- are formalized and machine-checked in Lean~4 against \texttt{mathlib}; the proof source accompanies this manuscript.

Thus the long wall remains a useful diagnostic.  It does not reveal a relative-phase obstruction in the restricted two-boundary model, but a complete time-symmetric theory still owes the dynamical results listed above.

\section{Conclusion}

The forward-only part of the long-wall objection is straightforward: if the empty-channel amplitude ratio tends to zero, the coherent field visibility tends to zero.  This is the behavior predicted for source-driven walking-droplet models under the undriven-branch assumptions.

For the scalar paraxial point-slit model considered here, a field conditioned on a candidate detector event has the standard two-slit relative phase,
\begin{equation}
 \Delta\phi=-\frac{kaY}{d}
 \quad(\mathrm{mod}\ 2\pi).
\end{equation}
If the two backward branches subsequently undergo identical nonzero scalar transfer, that propagation changes neither their amplitude ratio nor their relative phase.  This is a common-mode statement, not a solution of a material-wall boundary problem.

The calculation therefore removes one possible obstruction: extending a matched common channel need not introduce differential attenuation or an additional relative phase between the two future-conditioned branches.  It does not establish the strength of a self-consistent absorber response, the visibility of a field containing additional contributions, or the probability distribution of detector events.  It addresses one propagation component of the long-wall objection rather than settling the complete dynamical objection.

\begin{thebibliography}{99}

\bibitem{Andersen2015}
A.~Andersen, J.~Madsen, C.~Reichelt, S.~R.~Ahl, B.~Lautrup, C.~Ellegaard, M.~T.~Levinsen, and T.~Bohr,
``Double-slit experiment with single wave-driven particles and its relation to quantum mechanics,''
\emph{Physical Review E} \textbf{92}, 013006 (2015).

\bibitem{Bohr2016}
T.~Bohr, A.~P.~Andersen, and B.~Lautrup,
``Bouncing droplets, pilot-waves, and quantum mechanics,''
in \emph{Recent Advances in Fluid Dynamics with Environmental Applications}, pp.~335--349, Springer (2016).

\bibitem{Daniel2022}
A.~Daniel, B.~Ghosh, B.~Gorzkowski, and R.~Lapkiewicz,
``Demonstrating backflow in classical two beams' interference,''
\emph{New Journal of Physics} \textbf{24}, 123011 (2022).

\bibitem{WheelerFeynman1945}
J.~A.~Wheeler and R.~P.~Feynman,
``Interaction with the absorber as the mechanism of radiation,''
\emph{Reviews of Modern Physics} \textbf{17}, 157--181 (1945).

\bibitem{WheelerFeynman1949}
J.~A.~Wheeler and R.~P.~Feynman,
``Classical electrodynamics in terms of direct interparticle action,''
\emph{Reviews of Modern Physics} \textbf{21}, 425--433 (1949).

\bibitem{Cramer1986}
J.~G.~Cramer,
``The transactional interpretation of quantum mechanics,''
\emph{Reviews of Modern Physics} \textbf{58}, 647--687 (1986).

\bibitem{Wharton2007}
K.~B.~Wharton,
``Time-symmetric quantum mechanics,''
\emph{Foundations of Physics} \textbf{37}, 159--168 (2007).

\bibitem{Hasselmann1997}
K.~Hasselmann,
``The metron model: elements of a unified deterministic theory of fields and particles. Part 3: Quantum phenomena,''
\emph{Physics Essays} \textbf{10}, 64--86 (1997).

\bibitem{Hasselmann1998}
K.~Hasselmann,
``The metron model: towards a unified deterministic theory of fields and particles,''
arXiv:hep-th/9810086 (1998).

\bibitem{Hasselmann2004}
K.~Hasselmann and S.~Hasselmann,
``The metron model: a unified deterministic theory of fields and particles---a progress report,''
\emph{Proceedings of the Institute of Mathematics of the National Academy of Sciences of Ukraine} \textbf{50}, 788--795 (2004).

\end{thebibliography}

\end{document}
